
\section{Introduction}

Previous research has shown that large language models (LLMs) can be vulnerable to attacks at both training and inference time. Some examples of adversarial methods that have successfully elicited improper LLM generations include data poisoning \citep{du2022ppt} and adversarially generated prompts \citep{zou2023universal_adv_attacks}. Such attacks aim to manipulate the model and cause it to generate unsafe outputs that will be harmful for users. Improving language models' robustness to such adversarial prompts is therefore a key aspect of making such models safe for real-world deployment.

Multi-agent debate is a technique where multiple instances of a language model critique each others' responses, which can be seen as an extension of previous work such as chain-of-thought prompting and self-refinement. Previous work has used this to improve language models' factuality and reasoning, as well as their performance on downstream tasks such as evaluation. It seems reasonable that some amount of self-discussion could allow a language model to recognize potential downstream harms caused by its generation and revise its response. However, if an LLM ``agent'' outputs a toxic response, it could also poison the responses of other agents in the debate, leading to a more toxic output or no significant improvement in model toxicity. For this reason, while multi-agent debate has proven successful in other applications, it is important to study the debate dynamics specifically amongst LLMs prompted adversarially. In this way, we can not only analyze its efficacy in improving models' adversarial robustness, but also validate that this new technique does not introduce new risks at inference time.

In this work, we implement multi-agent debate between LLMs in the Llama-2 and GPT-3 families, using prompt engineering to simulate the effects of having poisoned models participate in a multi-agent debate. We then probe models' responses to prompts known to elicit toxic responses in single-agent, self-refine, and multi-agent debate settings. We find that models equipped with multi-agent debate at inference time generally produce less toxic responses to adversarial prompts, even for models that have previously been fine-tuned with methods such as reinforcement learning from human feedback. We also find that multi-agent debate outperforms methods like self-refinement in cases where an initial model is poisoned. This work motivates future exploration into multi-agent debate dynamics between LLMs and whether they can improve robustness in more realistic environments.
